\section{Construction of intertwining operators}\label{section5}

Let $G$ be a one-parameter metabelian nilpotent  Lie group  of
dimension $n$ with Lie algebra~$\frak g$.
 Let $\Gamma$ be a discrete uniform
subgroup of~$G$. Let $M = \exp(\frak m)$, where $\frak m$ is  a
rational abelian ideal in~$\frak g$ of codimension one.
 Let  $\bbase = (X_1, \ldots, X_n)$ be a
strong Malcev basis of $\frak g$ strongly based on~$\Gamma$ passing
through $[\frak g, \frak g]$ and $\frak m$. We put
\begin{gather}\label{alg-derivee}
[\frak g, \frak g] = \spann\lbrace X_1,\ldots, X_p\rbrace
\end{gather}
and
\[
\frak m
=\spann\set{X_1, \ldots, X_{n-1}}.
\]
  Let
\[
\nongen = \mathds{Z} X_{p+1}^*+\cdots+\mathds{Z}X_{n}^*
\]
and
\[
\gen= \{ l \in
\mathds{Z}X_1^*+\cdots +\mathds{Z}X_{n-1}^*:\ l\vert_{[\frak g,
\frak g]}\ne 0\}.
\]
 In the following, for $l \in \frak m^* \subset \frak g^*$ and $g
\in G$ we denote
\[
 \mathrm{Ad}_0^*g l =(\mathrm{Ad}^* g l)\vert_{\frak m}.
 \]


\begin{lemma}
 The subset $\gen$ is $  \mathrm{Ad}_0^* \Gamma$ invariant.
\end{lemma}
\begin{proof}
Let $\exp(\gamma) \in \Gamma$ and $l \in \gen$. Let $i=1, \ldots,
n-1$. By def\/inition of the coadjoint representation we have
\[
\langle \exp(-\gamma)\cdot l, X_i\rangle = \langle l, e^{\ad
\gamma}(X_i)\rangle.
\]
 On the other hand, since $M$ is normal in
$G$,  we have  that
\[
\exp\big(e^{\ad
\gamma}(X)\big) = \exp(\gamma) \exp{X_i} \exp(-\gamma) \in \Gamma \cap
M.
\] Since $M$ is abelian and
\[
\Gamma \cap M=\exp(\dsz X_1)\cdots
\exp(\dsz X_{n-1})
\] then
\[
\log(\Gamma) \cap \frak m=\dsz X_1\oplus\cdots\oplus \dsz
X_{n-1}.
\] It follows that  $\langle l, e^{\ad \gamma}(X_i)\rangle
\in \dsz$. Then
\begin{gather}\label{cond1}
\exp(-\gamma)\cdot l\in
\dsz X_1^*\oplus\cdots\oplus \dsz X_{n-1}^*.
\end{gather}
It remains to prove that
\begin{gather}\label{cond2}
(\exp(-\gamma)\cdot l)\vert_{[\frak g, \frak g]}\ne 0.
\end{gather}
We have
\[
\frak g(\exp(-\gamma)\cdot l)=\Ad\exp(-\gamma)(\frak g(l)).
\]
Since $\frak g(l)\ne \frak g$, then $\frak g(\exp(-\gamma)\cdot
l)\ne \frak g$ and therefore (\ref{cond2}) holds. Consequently, from
(\ref{cond1}) and~(\ref{cond2}) we have $\exp(-\gamma)\cdot l \in
\gen$.
\end{proof}



Let $\mathcal{C}(G/\Gamma)$ be the space of all complex valued
continuous functions $\xi$ on $G$ satisfying
\begin{gather}\label{periode}
\xi(g\gamma)= \xi(g), \end{gather}
 for any $g$ in $G$ and $\gamma$
in $\Gamma$.




\begin{lemma}\label{periode2} Let $l\in \nongen$, we have
\begin{itemize}\itemsep=0pt
  \item[$(1)$] $\chi_l\vert_{\Gamma}=1$.
  \item[$(2)$] For $\xi \in \mathcal{C}(G/\Gamma)$, and $g\in G$, the function
\[
G/\Gamma\rightarrow \dsc,\qquad m\Gamma\mapsto \xi(gm)
\chi_l(m)
\]
is well defined.
\end{itemize}
\end{lemma}

\begin{proof} $(1)$ Let $\gamma\in \Gamma$, then there exist
$t_1, \ldots, t_n\in \dsz$ such that
\[
\gamma=\exp(t_1X_1)\cdots \exp(t_nX_n).
\]
By the Baker--Campbell--Hausdorf\/f formula we have
\[
\log(\gamma)\equiv t_{p+1}X_{p+1}+\cdots+t_n X_n\pmod{[\frak g,
\frak g]}.
\]
 Since $l\vert_{[\frak g, \frak g]}=0$ then $\chi_l(\gamma)=1$.


$(2)$ Let $\gamma\in \Gamma$. We have
\begin{gather*}
 \xi(gm\gamma)
\chi_l(m\gamma)
   = \xi(gm)
\chi_l(m\gamma) \overset{\text{by (\ref{periode})}}{=}    \xi(gm) \chi_l(m) \chi_l(\gamma)
 = \xi(gm) \chi_l(m).
\end{gather*}
 This completes the proof of the lemma.
\end{proof}

Let $\Sigma$ be a crosssection for $\Gamma$-orbits in $\gen$. Let
\begin{gather*}
 \rho =   \bigoplus_{l \in \Sigma}\ind_M^G\chi_l +
\bigoplus_{l\in \nongen} \chi_l.
\end{gather*}


We are now in a position  to formulate the following

\begin{theorem}\label{des-operator-special}
The operator  $\U$ defined for all  $\xi\in \mathcal{C}(G/\Gamma)$
and $g\in G$ by
 \begin{gather*}%\label{operator-special}
 \U(\xi)(l)(g) = \begin{cases}
\displaystyle     \int_{M/M\cap \Gamma} \xi(gm) \chi_l(m) d\dot{m} & \mbox{if }  l\in \Sigma,\vspace{1mm}\\
\displaystyle     \int_{G/ \Gamma} \xi(gm) \chi_l(m) d\dot{m} & \mbox{if }  l\in \nongen
    \end{cases}
\end{gather*}
 is an isometric linear operator having value in the Hilbert space
 $\Hi_\rho$ of $\rho$  and can be extended on
  $\mathbf{L}^2(G/\Gamma)$ to an
intertwining operator of $\tau$ and $\rho$.
\end{theorem}

\begin{proof} Clearly
for $ l \in \Sigma\cup \nongen$, the function $ \U(\xi)(l)$
satisf\/ies the covariance relation (\ref{covarience}). First, we
establish that $\U$ is well def\/ined and isometric. Let $\xi\in
\mathcal{C}(G/\Gamma)$. Then
\begin{gather*}
  \|\U(\xi)\|^2  = \sum_{l \in \Sigma}\|\U(\xi)(l)\|^2_{L^2(G/M, l)} + \sum_{l\in
  \nongen}
    \|\U(\xi)(l)\|^2.
\end{gather*}
We proceed to calculate the f\/irst sum. For $x=(x_1,\ldots, x_{n-1})
\in \mathds{R}^{n-1}$ and $ t\in \mathds{R}$, let
\begin{gather*}
\delta(t, x) = \mathbf{E}_{\bbase}((x, t))= \exp(t X_n) \exp(x_{n-1} X_{n-1})\cdots \exp(x_1
X_1),
\\
\sum_{l\in \Sigma} \|\U(\xi)(l)\|^2_{L^2(G/M, l)}
   = \sum_{l\in \Sigma} \int_{G/M} | \U(\xi)(l)(g)|^2     d\dot{g }\\
 \hphantom{\sum_{l\in \Sigma} \|\U(\xi)(l)\|^2_{L^2(G/M, l)}} {}
  =  \sum_{l\in \Sigma} \int_{G/M} \bigg| \int_{M/M\cap \Gamma}
\xi(gm) \chi_l(m) d\dot{m}\bigg|^2     d\dot{g}\\
\hphantom{\sum_{l\in \Sigma} \|\U(\xi)(l)\|^2_{L^2(G/M, l)}} {} = \sum_{l\in \Sigma} \sum_{s\in \mathds{Z}} \int_{\tore} \bigg|
\int_{M/M\cap \Gamma} \xi(\delta(t, 0)\delta(s, 0)  m) \chi_l (m)
 d\dot{m} \bigg|^2  dt \\
\hphantom{\sum_{l\in \Sigma} \|\U(\xi)(l)\|^2_{L^2(G/M, l)}} {} =
\sum_{l\in \Sigma} \sum_{s\in \mathds{Z}} \int_{\tore} \bigg|
\int_{M/M\cap \Gamma} \xi(\delta(t, 0)  m) \chi_l ( \delta(s,
0)^{-1} m \delta(s, 0))
 d\dot{m} \bigg|^2  dt\\
\hphantom{\sum_{l\in \Sigma} \|\U(\xi)(l)\|^2_{L^2(G/M, l)}} {}  =  \sum_{l\in \Sigma} \sum_{s\in
\mathds{Z}} \int_{\tore}  \bigg| \int_{M/M\cap \Gamma} \xi(\delta(t, 0)
m) \chi_{\Ad^*(\delta(s, 0))l} (m)
 d\dot{m} \bigg|^2  dt .
\end{gather*}
As the mapping $\Sigma\times \mathds{Z}\rightarrow \gen$, $(l, s) \mapsto \Ad^*(\delta(s, 0))l$
 is bijective, then
 \begin{gather*}
 \sum_{l\in \Sigma} \|\U(\xi)(l)\|^2_{L^2(G/M, l)}
   =  \sum_{l\in \gen}  \int_{\tore}  \bigg| \int_{M/M\cap
\Gamma} \xi(\delta(t, 0) m) \chi_l (m)
 d\dot{m} \bigg|^2  dt  \\
 \hphantom{\sum_{l\in \Sigma} \|\U(\xi)(l)\|^2_{L^2(G/M, l)}}{}
  =  \sum_{l\in \gen}  \int_{\tore}
 \bigg|
 \int_{\tore^{n-1}} \xi(\delta(t, x))
  \chi_l (\delta(0, x))
 dx  \bigg|^2  dt.
\end{gather*}
Next, we compute
\begin{gather*}
\sum_{l\in
  \nongen}
    \|\U(\xi)(l)\|^2  =
\sum_{l\in \nongen}
    | \U(\xi)(l)(e) |^2  = \sum_{l\in \nongen}\bigg|\int_{G/\Gamma}
    \xi(g) \chi_l(g) d\dot{g} \bigg|^2\\
\hphantom{\sum_{l\in   \nongen}     \|\U(\xi)(l)\|^2}
 = \sum_{l\in \nongen}\bigg|\int_{\tore^{n}}
    \xi(\delta(t, x)) \chi_l(\delta(t, x)
 ) dt dx \bigg|^2\\
\hphantom{\sum_{l\in   \nongen} \|\U(\xi)(l)\|^2}
 = \sum_{l_1\in \nongen_1} \sum_{l_2\in \nongen_2} \bigg|\int_{\tore^{n}}
    \xi(\delta(t, x)) \chi_{l_1}(\delta(0, x))
     \chi_{l_2}(\delta(t, 0) ) dt dx \bigg|^2\\
\hphantom{\sum_{l\in   \nongen} \|\U(\xi)(l)\|^2}
 (\text{where} \ \nongen_1 = \mathds{Z} X_{p+1}^*+\cdots+\mathds{Z}X_{n-1}^*,
     \nongen_2 =\mathds{Z}X_{n}^* \text{ and } l = l_1+l_2
      \in \nongen_1\oplus \nongen_2) \\
\hphantom{\sum_{l\in   \nongen} \|\U(\xi)(l)\|^2}
 = \sum_{l_1\in \nongen_1} \int_{\tore} \bigg|\int_{\tore^{n-1}}
    \xi(\delta(t, x))
 \chi_{l_1}(\delta(0, x))  dx \bigg|^2 dt,
\end{gather*}
where we have applied Parseval's equality with respect to the
variable $t$. Summarizing, we have
\begin{gather*}
  \|\U(\xi)\|^2    =  \sum_{l\in \gen}  \int_{\tore}
 \bigg|
 \int_{\tore^{n-1}} \xi(\delta(t, x))
  \chi_l (\delta(0, x))
 dx  \bigg|^2  dt \\
\hphantom{\|\U(\xi)\|^2    =}{}
+  \sum_{l\in \nongen_1} \int_{\tore}
\bigg|\int_{\tore^{n-1}}
    \xi(\delta(t, x))\chi_{l}(\delta(0, x))
dx \bigg|^2 dt.
\end{gather*}

It is clear that $\gen$ and $\nongen_1$ are disjoint. In fact,
suppose that $\gen\cap\nongen_1\ne \varnothing$. Let $l\in
\gen\cap\nongen_1$. The condition $l\in \nongen_1$ implies that
$l\vert_{[\frak g, \frak g]}= 0$ since $[\frak g, \frak g]=\dsr
X_1\oplus\cdots\oplus \dsr X_p$ (see~(\ref{alg-derivee})). This
contradicts the fact that $l\vert_{[\frak g, \frak g]}\ne 0$ because
$l\in \gen$. Now, let
\begin{gather}\label{disjoint-reunion}
\frak m_{\mathds{Z}}^* =\gen\sqcup \nongen_1.
\end{gather}
 It is
clear that
\[
\frak m_{\mathds{Z}}^* = \mathds{Z}X_1^*+\cdots
+\mathds{Z}X_{n-1}^*.
\]
 Consequently, we obtain
\begin{gather*} \|\U(\xi)\|^2
= \int_{\tore} \sum_{l\in \frak
m_{\mathds{Z}}^*} \left|\int_{\tore^{n-1}}
    \xi(\delta(t, x)) \chi_{l}(\delta(0, x))
dx \right|^2 dt   = \int_{\tore} \int_{\tore^{n-1}}
    |\xi(\delta(t, x))
 |^2 dx  dt \\
 \hphantom{\|\U(\xi)\|^2 }{}
 \overset{(\text{by Parseval's equality})}{=}  \int_{\tore^{n}}
    |\xi(\delta(t, x))
 |^2 dx  dt
   = \|\xi\|_{\mathscr{H}_\tau}^2.
\end{gather*}
The operator $\U$ being now isometric, it can be extended to an
isometry of $\mathbf{L}^2(G/\Gamma)$. It remains to verify that $\U$
is an intertwining operator for $\tau$ and $\rho$. It suf\/f\/ices then
to prove that  $\U\circ \tau(g)(\xi) = \rho(g) \circ \U(\xi)$ for
every  $g$ in $G$ and $\xi$ in $\mathcal{C}(G/\Gamma)$. Let $l\in
\Sigma, a\in G$. We compute
\[
\U\circ \tau(g)(\xi)(l)(a)=\int_{M/M\cap \Gamma} \tau(g)(\xi)(am)
\chi_l(m) d\dot{m}= \int_{M/M\cap \Gamma} \xi(g^{-1}am) \chi_l(m)
d\dot{m}.
\]
 On the other hand
\begin{gather*}
\rho(g) \circ \U(\xi)(l)(a)=\ind_M^G\chi_l(g)(\U(\xi)(l))(a)=
\U(\xi)(l)\big(g^{-1}a\big)\\
\hphantom{\rho(g) \circ \U(\xi)(l)(a)}{}
= \int_{M/M\cap \Gamma} \xi(g^{-1}am) \chi_l(m)
d\dot{m}.
\end{gather*}
Thus
\[
\U\circ \tau(g)(\xi)(l)(a)=\rho(g) \circ \U(\xi)(l)(a).
\]
Similarly, we prove that the same equality holds if $l\in \nongen$.
 Consequently,
 the following
diagram:
 \begin{gather}\label{intertwining}
\begin{CD} \mathscr{H}_\tau@>{\tau(g)}>> \mathscr{H}_\tau\\
@V{\U}VV  @V{\U}VV\\
\mathscr{H}_\rho @> {\rho(g)}>> \mathscr{H}_\rho
\end{CD}
\end{gather}
is commutative.
\end{proof}








Next, we show that $\U$ is an invertible operator. For this, let
\[
 \Hi_\rho^c =  \bigoplus_{l\in \Sigma} \mathcal{C}_c(G/M,
\chi_l) \oplus \bigoplus_{l\in \nongen} \mathds{C} \chi_l \subset
\mathscr{H}_\rho,
\]
 where $\mathcal{C}_c(G/M, \chi_l)$ is the space
of complex valued continuous functions $\xi$ on $G$ satisfying
\begin{gather}\label{covariance2}
\xi(gm)=\chi_l^{-1}(m) \xi(g) \qquad
  \forall\, g\in G, \ \ m\in M,
\end{gather}
and having a compact support modulo $M$.



\begin{lemma} Let $K\in \Hi_\rho^c$, $l\in \Sigma$ and $g\in G$. The
function
\[
\Gamma/\Gamma\cap
M \rightarrow \dsc, \qquad \gamma (\Gamma\cap M)\mapsto
K(l)(g\gamma)
\]
 is well defined.
\end{lemma}

\begin{proof} Let $\gamma\in \Gamma$ and $\gamma'\in\Gamma\cap M$.
Since $K(l)\in \mathcal{C}_c(G/M, \chi_l)$ then $K(l)$ satisf\/ies the
covariance relation (\ref{covariance2}). Hence $
K(l)(g\gamma\gamma') = \chi_l^{-1}(\gamma') K(l)(g\gamma)$. On the
other hand, as $\Gamma\cap M= \exp(\dsz X_1)\cdots \exp(\dsz
X_{n-1})$ and $M$ is abelian, then $\log(\Gamma)\cap \frak m= \dsz
X_1\oplus \dots \oplus \dsz X_{n-1}$. It follows that
$\chi_l\vert_{\Gamma\cap M}=1$, in particular, $\chi_l(\gamma')=1$.
Then $ K(l)(g\gamma\gamma') = K(l)(g\gamma)$.
\end{proof}

Let
\begin{gather*}
 \Hi_\rho^0 =\{K\in \Hi_\rho^c:\ K(l)  \mbox{is a zero function everywhere, except
for f\/inite number of  } l\}.
\end{gather*}
The space $ \Hi_\rho^0$ is dense in $\Hi_\rho$.


\begin{lemma} For $K\in \Hi_\rho^0 $ and $g\in G$, the sum
\begin{gather}\label{inverse}
  \sum_{l\in \Sigma} \sum_{\gamma\in \Gamma/\Gamma\cap
M} K(l)(g\gamma)+ \sum _{l\in \nongen} K(l)(g), \qquad g\in
G,
\end{gather} has only finitely many nonzero terms.
\end{lemma}

\begin{proof} Let  $l\in \Sigma\cup \nongen$.
For $g\in G$, let
\[
S_g=\set{\gamma\in \Gamma:\ g\gamma\in \mathrm{supp}(K(l))}.
\]
Then there is an integer $n_{K(l)}$ independent of $g\in G$ such
that $S_g$ is the union of at most $n_{K(l)}$ cosets of~$\Gamma\cap M$ \cite[Lemma~3.2]{Cor4}. This observation shows that the sum
over $\Gamma/\Gamma\cap M$  in~(\ref{inverse}) has at most
$n_{K(l)}$ nonzero entries for each~$g\in G$. Consequently, the sum~(\ref{inverse}) is a sum of f\/inite terms.
\end{proof}


We def\/ine the  operator $\V$  on the dense subspace $ \Hi_\rho^0 $
by
\begin{gather*}
\V(K)(g) =  \sum_{l\in \Sigma} \sum_{\gamma\in \Gamma/\Gamma\cap
M} K(l)(g\gamma)+ \sum _{l\in \nongen} K(l)(g),\qquad g\in G.
\end{gather*}

